% ============================================================
%  Příprava na maturitu – Mechatronika
%  Kompiluj: pdflatex maturita_mechatronika.tex
%  Kódování: UTF-8
% ============================================================

\documentclass[a4paper,10pt]{article}

% --- Balíčky ---
\usepackage[T1]{fontenc}
\usepackage[utf8]{inputenc}
\usepackage[czech]{babel}
\usepackage[left=2.4cm, right=1.8cm, top=1.8cm, bottom=1.8cm]{geometry}
\usepackage{lmodern}
\usepackage{microtype}
\usepackage{parskip}
\usepackage{titlesec}
\usepackage{enumitem}
\usepackage{xcolor}
\usepackage{hyperref}
\usepackage{tabularx}
\usepackage{booktabs}
\usepackage{array}
\usepackage{tikz}
\usepackage{eso-pic}
\usepackage{mdframed}
\usepackage{amsmath}
\usepackage{amssymb}
\usepackage{pgfplots}
\pgfplotsset{compat=1.18}
\usetikzlibrary{arrows.meta, positioning, shapes.geometric, calc}

% --- Barvy ---
\definecolor{accent}{HTML}{03254E}
\definecolor{stripeblue}{HTML}{568EA3}
\definecolor{medgray}{HTML}{888888}
\definecolor{lightblue}{HTML}{EAF2F8}
\definecolor{lightyellow}{HTML}{FDFBE4}
\definecolor{lightgreen}{HTML}{EAF7EA}

% --- Modrý pruh vlevo ---
\AddToShipoutPictureBG{%
  \begin{tikzpicture}[remember picture, overlay]
    \fill[stripeblue]
      (current page.north west)
      rectangle
      ([xshift=10mm] current page.south west);
    \fill[accent!60]
      ([xshift=10mm] current page.north west)
      rectangle
      ([xshift=12mm] current page.south west);
  \end{tikzpicture}%
}

\hypersetup{colorlinks=true, urlcolor=accent, linkcolor=accent}

\titleformat{\section}
  {\large\bfseries\color{accent}}
  {\thesection.}
  {0.5em}
  {}
  [\color{accent}\titlerule]

\titlespacing{\section}{0pt}{16pt}{6pt}

\newmdenv[
  backgroundcolor=lightblue,
  linecolor=stripeblue,
  linewidth=1pt,
  roundcorner=4pt,
  innerleftmargin=8pt,
  innerrightmargin=8pt,
  innertopmargin=6pt,
  innerbottommargin=6pt,
  skipabove=4pt,
  skipbelow=4pt
]{pojmybox}

\newmdenv[
  backgroundcolor=lightyellow,
  linecolor=accent!40,
  linewidth=1pt,
  roundcorner=4pt,
  innerleftmargin=8pt,
  innerrightmargin=8pt,
  innertopmargin=6pt,
  innerbottommargin=6pt,
  skipabove=4pt,
  skipbelow=8pt
]{vysvetlenibox}

\newmdenv[
  backgroundcolor=lightgreen,
  linecolor=accent!30,
  linewidth=1pt,
  roundcorner=4pt,
  innerleftmargin=8pt,
  innerrightmargin=8pt,
  innertopmargin=6pt,
  innerbottommargin=6pt,
  skipabove=4pt,
  skipbelow=8pt
]{vzorcebox}

\pagestyle{plain}

% ============================================================
\begin{document}

% --- Záhlaví ---
\begin{minipage}[t]{0.75\textwidth}
    {\Huge\bfseries\color{accent} Příprava na maturitu}\\[4pt]
    {\large\color{stripeblue} Mechatronika}\\[2pt]
    \textcolor{medgray}{\small Suchánek Lukáš \quad SPŠ Prosek \quad 2026}
\end{minipage}
\hfill
\begin{minipage}[t]{0.22\textwidth}
    \raggedleft
    \begin{tikzpicture}
        \draw[color=stripeblue, rounded corners=4pt, line width=1pt]
            (0,0) rectangle (3,1.2)
            node[midway, align=center, color=accent, font=\small\bfseries]
            {Otázek: 22};
    \end{tikzpicture}
\end{minipage}

\vspace{6pt}
\noindent\color{accent}\rule{\linewidth}{1.5pt}
\vspace{2pt}

% ============================================================
\section{Mechatronika, mechatronický návrh, systém a výrobek}

\begin{pojmybox}
\textbf{Klíčové pojmy:}
mechatronika, mechanika, elektronika, informatika,
senzor, aktuátor, řídící systém, systémový přístup, mechatronický výrobek,
V-model návrhu, SiL, HiL, zpětná vazba
\end{pojmybox}

\begin{vysvetlenibox}
\textbf{Definice a vznik oboru}\\
Mechatronika je interdisciplinární obor kombinující \textbf{mechaniku},
\textbf{elektroniku} a \textbf{informatiku} za účelem návrhu inteligentních
systémů a výrobků. Termín poprvé použila japonská firma Yaskawa Electric
v roce 1969. Dnes je mechatronika základem moderní průmyslové automatizace,
automobilového průmyslu, robotiky i spotřební elektroniky.

\textbf{Historický vývoj mechatroniky}\\
Vývoj mechatroniky lze rozdělit do několika fází:
\begin{itemize}[noitemsep, leftmargin=1.4em, topsep=4pt]
    \item \textbf{1969--1980} -- vznik termínu, první mechatronické systémy (přádací stroje, CNC)
    \item \textbf{1980--1995} -- rozvoj mikroprocesorů, masové nasazení v automobilkách
    \item \textbf{1995--2010} -- nástup průmyslových sběrnic, decentralizovaná řízení
    \item \textbf{2010--současnost} -- Industry 4.0, IoT, umělá inteligence, kyberfyzikální systémy
\end{itemize}

\textbf{Složky mechatronického systému}\\
Každý mechatronický výrobek obsahuje tři neoddělitelné složky:
\begin{itemize}[noitemsep, leftmargin=1.4em, topsep=4pt]
    \item \textbf{Mechanická část} -- nosná konstrukce, převody, kinematické
          vazby; zajišťuje pohyb a přenos sil
    \item \textbf{Elektronická část} -- senzory (měří fyzikální veličiny),
          aktuátory (vykonávají pohyb), výkonová elektronika (řídí výkon)
    \item \textbf{Řídicí část} -- mikrokontrolér nebo PLC s algoritmem;
          zpracovává data ze senzorů a vydává příkazy aktuátorům
\end{itemize}

\textbf{Informační toky v mechatronickém systému}\\
Signál putuje systémem v uzavřené smyčce:
\begin{enumerate}[noitemsep, leftmargin=1.4em, topsep=4pt]
    \item Fyzikální veličina (teplota, tlak, poloha) působí na \textbf{senzor}
    \item Senzor převede veličinu na elektrický signál (napětí, proud, frekvence)
    \item Signál je kondicionován (zesílen, filtrován) a převeden A/D převodníkem
    \item \textbf{Řídicí jednotka} vyhodnotí stav a vypočítá akční zásah
    \item Akční člen (\textbf{aktuátor}) provede pohyb nebo změnu výkonu
    \item Nový stav je opět měřen senzorem -- smyčka se uzavírá
\end{enumerate}

\textbf{Systémový přístup a V-model}\\
Na rozdíl od tradičního sekvenčního návrhu (nejprve mechanika, pak elektronika,
nakonec software) systémový přístup navrhuje všechny složky \textit{současně}.
Standardním postupem je \textbf{V-model}:

\vspace{4pt}
\begin{center}
\begin{tikzpicture}[scale=0.85,
  box/.style={rectangle, draw=accent, fill=lightblue, rounded corners=3pt,
              minimum width=2.8cm, minimum height=0.6cm, font=\small,
              text centered}]
  \node[box] (spec)    at (0,0)    {Specifikace požadavků};
  \node[box] (arch)    at (2,-1)   {Architektura systému};
  \node[box] (detail)  at (4,-2)   {Detailní návrh};
  \node[box] (impl)    at (6,-3)   {Implementace};
  \node[box] (unit)    at (8,-2)   {Unit testování};
  \node[box] (integ)   at (10,-1)  {Integrační testy};
  \node[box] (valid)   at (12,0)   {Validace systému};
  \draw[>-, accent, thick]
    (spec) -- (arch) -- (detail) -- (impl) -- (unit) -- (integ) -- (valid);
  \draw[dashed, medgray]
    (spec.south east) -- (valid.south west)
    node[midway, below, font=\scriptsize, color=medgray] {čas / fáze};
\end{tikzpicture}
\end{center}
\end{vysvetlenibox}

\begin{vysvetlenibox}
\textbf{Fáze V-modelu v praxi:}
\begin{itemize}[noitemsep, leftmargin=1.4em, topsep=4pt]
    \item \textbf{Specifikace požadavků} -- co má systém umět (rychlost, přesnost, bezpečnost)
    \item \textbf{Architektura systému} -- rozdělení na mechaniku, elektroniku, software
    \item \textbf{Detailní návrh} -- CAD modely, schémata elektroniky, algoritmy
    \item \textbf{Implementace} -- výroba prototypu, programování, osazování
    \item \textbf{Unit testování} -- testování jednotlivých komponent
    \item \textbf{Integrační testy} -- ověření spolupráce všech složek
    \item \textbf{Validace systému} -- splnění původních požadavků zákazníka
\end{itemize}

\textbf{SiL a HiL testování}\\
Moderní vývoj bez fyzického prototypu v raných fázích:
\begin{itemize}[noitemsep, leftmargin=1.4em, topsep=4pt]
    \item \textbf{SiL} (Software in the Loop) -- software řídicího algoritmu
          běží na PC a simuluje chování soustavy; testuje logiku bez HW.
          \textit{Příklad:} Simulace PID regulátoru v MATLAB/Simulink.
    \item \textbf{HiL} (Hardware in the Loop) -- reálný řídicí HW (PLC, MCU)
          je připojen k simulátoru fyzikální soustavy; testuje se odezva na
          reálná vstupní data v reálném čase.
          \textit{Příklad:} Ověření airbag řídicí jednotky s simulátorem nárazu.
    \item \textbf{MiL} (Model in the Loop) -- testování matematického modelu
          proti požadavkům; nejranější fáze verifikace.
\end{itemize}

\textbf{Příklady mechatronických výrobků}

\begin{center}
\begin{tabularx}{0.95\linewidth}{@{} l X X X @{}}
    \toprule
    \textbf{Výrobek} & \textbf{Mechanika} & \textbf{Elektronika}
                     & \textbf{Řídicí systém} \\
    \midrule
    ABS brzdy        & brzdový třmen, pístnice & senzory otáček kol
                     & ECU, PID regulátor \\
    Průmyslový robot & kinematický řetězec, klouby & servomotory, enkodéry
                     & trajektoriový kontrolér \\
    Pračka           & buben, ložiska, čerpadlo & motor, ventily, senzory
                     & PLC / MCU, program cyklů \\
    CNC frézka       & portál, vřeteno, osy & servopohony, čidla polohy
                     & CNC řídicí systém, G-kód \\
    \bottomrule
\end{tabularx}
\end{center}

\textbf{Detailní příklad: Antiblokovací brzdový systém (ABS)}\\
ABS je ukázkový mechatronický systém v automobilu:
\begin{itemize}[noitemsep, leftmargin=1.4em, topsep=4pt]
    \item \textbf{Mechanika:} Brzdový třmen tlačí brzdové destičky na kotouč.
          Pístnice přenáší sílu z brzdového pedálu.
    \item \textbf{Elektronika:} Indukční senzory u každého kola měří okamžitou
          úhlovou rychlost. Při rozdílu otáček kol detekují smyk.
    \item \textbf{Řízení:} Řídicí jednotka (ECU) vyhodnocuje signály ze senzorů
          100× za sekundu. Při detekci smyku otevře elektromagnetický ventil,
          který sníží tlak v brzdovém systému a kolo se opět roztočí.
\end{itemize}

\textbf{Výhody mechatronického přístupu:}
\begin{itemize}[noitemsep, leftmargin=1.4em, topsep=4pt]
    \item \textbf{Optimalizace celku} -- lze najít lepší řešení než optimalizací každé části zvlášť
    \item \textbf{Flexibilita} -- změny ve funkci se často řeší úpravou software
    \item \textbf{Diagnostika} -- systém umí detekovat poruchy a informovat obsluhu
    \item \textbf{Adaptivita} -- moderní systémy se učí a přizpůsobují podmínkám
\end{itemize}

\textbf{Nevýhody a výzvy:}
\begin{itemize}[noitemsep, leftmargin=1.4em, topsep=4pt]
    \item \textbf{Komplexita} -- těžší pochopení a ladění celého systému
    \item \textbf{Bezpečnost} -- software může obsahovat chyby s vážnými důsledky
    \item \textbf{Kyberbezpečnost} -- připojené systémy jsou zranitelné vůči útokům
    \item \textbf{Odborníci} -- potřeba inženýrů s přesahem do více oborů
\end{itemize}
\end{vysvetlenibox}

% ============================================================

% ============================================================
\end{document}
